% SECTION: `\documentclass`
\documentclass{article}

% IMPORTS: Packages
% Fonts
% \usepackage[utf8]{inputenc}   % Standardise font; only load one of these
\usepackage[T1]{fontenc}        % Standardise font; only load one of these
% AMS and math typesetting
\usepackage{amssymb}      % duh
\usepackage{amsmath}      % duh
\usepackage{amsthm}       % `\qed` etc
\usepackage{mathtools}    % idk
\usepackage{eucal}        % Nice `\mathcal`
% \usepackage{esint}        % Nice integrals
\usepackage{eqparbox}     % Among other things, `\eqmathbox`
% Graphics and environments
\usepackage{tabularx}                 % Extensible tables
\usepackage{graphicx}                 % Images
\usepackage{wrapfig}                  % Wrapping text around figures
\usepackage{subcaption}               % Subfigures
\usepackage[outputdir=build]{minted}  % Syntax highlighted code; NOTE: Compile with `--shell-escape`
\usepackage{enumitem}                 % Nice enumeration, including `\begin{itemize}[nosep]`
\usepackage{mdframed}                 % Framed environments
\usepackage{cprotect}                 % Protecting section headings
\usepackage{placeins}                 % \FloatBarrier
% Tikz
\usepackage{tikz}                     % duh
\usetikzlibrary{
  % automata,                           % FSMs, Turing Machines, etc.
  arrows,                             % Many arrows :D
  positioning,                        % Fine-grained positioning
  calc                                % Calculated positioning
}
\usepackage{tikz-among-us}            % duh
\usepackage{tikz-cd}                  % duh

% Arrows :D
\usepackage{extarrows}
% Landscape pages
\usepackage{pdflscape}
% Colours
\usepackage{xcolor}
% Citation and bibliography
\usepackage{bookmark}
\usepackage[
  backend = biber, 
  style = alphabetic,
  sorting = debug
]{biblatex}
\addbibresource{bibliography.bib}
% Page geometry layout (margins etc)
\usepackage{geometry}
\geometry{
  a4paper,
  total={150mm,237mm},
  left=30mm,
  top=30mm,
} % A4 dimensions: 210 x 297
% Hyperlinks
\usepackage{hyperref}   % NOTE: Load hyperref last!
\hypersetup{
  colorlinks = true,
  linkcolor = blue,
  urlcolor = cyan,
}



% SECTION: Invert page colour
% Uncomment to invert page colour:
% \pagecolor[rgb]{0.1, 0.1, 0.1}
% \color[rgb]{0.9, 0.9, 0.9}



% SECTION: Assignment-specific
\newcommand{\FIXME}[1]{\textcolor{red}{FIXME:} \textcolor{red}{#1}}

\newcommand{\q}[1]{\section*{Question #1} \stepcounter{section}}
\newcommand{\sq}[2]{\subsection*{Question #1#2} \stepcounter{subsection}}
\newcommand{\ssq}[3]{\subsubsection*{Question #1#2 part #3} \stepcounter{subsubsection}}

\newcommand{\setq}[3]{
  \setcounter{section}{#1}
  \setcounter{subsection}{#2}
  \setcounter{subsubsection}{#3}
}
\newcommand{\incq}{
  \stepcounter{section}
}
\newcommand{\incsq}{
  \stepcounter{subsection}
}
\newcommand{\incssq}{
  \stepcounter{subsubsection}
}

\newcommand{\tq}{
  \section*{
    Question \arabic{section}
  }
}
\newcommand{\tsq}{
  \section*{
    Question \arabic{section}\alph{subsection}
  }
}
\newcommand{\tssq}{
  \section*{
    Question \arabic{section}\alph{subsection} part \roman{subsubsection}
  }
}

\newmdenv[]{assq}



% SECTION: eqmathbox
% https://tex.stackexchange.com/a/34412/5764
\makeatletter
\NewDocumentCommand{\eqmathbox}{o O{c} m}{%
  \IfValueTF{#1}
    {\def\eqmathbox@##1##2{\eqmakebox[#1][#2]{$##1##2$}}}
    {\def\eqmathbox@##1##2{\eqmakebox{$##1##2$}}}
  \mathpalette\eqmathbox@{#3}
}
\makeatother

% SECTION: Convenient spacing
\newcommand{\skoop}{\vspace{\baselineskip}}
\newcommand{\pan}{\par\noindent}

% SECTION: Standard symbols
\newcommand*\df{\mathop{}\!\mathrm{d}}

\newcommand{\deriv}[1]{\frac{\df}{\df #1}}
\newcommand{\derivup}[2]{\frac{\df #1}{\df #2}}
\newcommand{\pderiv}[1]{\frac{\partial}{\partial #1}}
\newcommand{\pderivup}[2]{\frac{\partial #1}{\partial #2}}

\newcommand{\powset}[1]{\mathcal{P}\left(#1\right)}
\newcommand{\pre}{\leftarrow}

\newcommand{\one}{\mathbf{1}}
\newcommand{\zero}{\mathbf{0}}

% \restrict{f}_{A} ~~> \left. f \right\rvert_{A}
\newcommand{\restrict}[1]{\left. #1 \right\rvert}

\newcommand{\actson}{\mathrel{\curvearrowright}}

% SECTION: `\set`, my favourite <3
\DeclarePairedDelimiterX\set[1]\lbrace\rbrace{\def\given{\;\delimsize\vert\;}#1}

% SECTION: Matrix extension
\makeatletter
\renewcommand*\env@matrix[1][*\c@MaxMatrixCols c]{%
  \hskip -\arraycolsep
  \let\@ifnextchar\new@ifnextchar
  \array{#1}}
\makeatother

% SECTION: Math operators
% Set theory
\DeclareMathOperator{\card}{card}

% Category theory
\DeclareMathOperator{\Set}{\mathbf{Set}}
\DeclareMathOperator{\Cat}{\mathbf{Cat}}
\DeclareMathOperator{\Top}{\mathbf{Top}}
\DeclareMathOperator{\Grp}{\mathbf{Grp}}
\DeclareMathOperator{\Vect}{\mathbf{Vect}}
\DeclareMathOperator{\NormedSpace}{\mathbf{NS}}
\DeclareMathOperator{\id}{id}

% Linear algebra
\DeclareMathOperator{\spans}{span}
\DeclareMathOperator{\rank}{rank}
\DeclareMathOperator{\tr}{tr}

% Matrix-vector
\renewcommand{\vec}[1]{\mathbf{#1}}
\newcommand{\transpose}{{\top}}

% Group theory
\DeclareMathOperator{\im}{im}

\DeclareMathOperator{\GL}{GL}
\DeclareMathOperator{\SL}{SL}

\DeclareMathOperator{\Stab}{Stab}
\DeclareMathOperator{\Orb}{\mathcal{O}}

% Graph theory
\newcommand{\headless}{\relbar\mkern-9mu\relbar} % Headless arrow, used like $a \rightarrow b$

% Common functions
\DeclareMathOperator{\sgn}{sgn}

% Real analysis

% Complex analysis
\let\Re\relax\DeclareMathOperator{\Re}{Re}
\let\Im\relax\DeclareMathOperator{\Im}{Im}
\DeclareMathOperator{\Arg}{Arg}

% Metric space topology
\DeclareMathOperator{\diam}{diam}

% Vector calculus
\DeclareMathOperator{\divs}{div}
\renewcommand{\div}{\divs}
\DeclareMathOperator{\curl}{curl}

% Differential geomerty
% \DeclareMathOperator{\fff}{\RN{1}}
% \DeclareMathOperator{\sff}{\RN{2}}
% \newcommand{\realproj}{\mathbb{R}\mathbb{P}}
% \newcommand{\flatmetric}{\bar{g}}
% \newcommand{\roundmetric}{\overset{\circ}{g}}
% \newcommand{\realprojmetric}{\hat{g}}
% \DeclareMathOperator{\smvfs}{\mathfrak{X}}
% \newcommand{\connect}{\nabla}

% Probability theory
% \DeclareMathOperator{\PP}{\mathbb{P}}
% \DeclareMathOperator{\EE}{\mathbb{E}}
% \DeclareMathOperator{\Var}{\mathbb{V}ar}
% \DeclareMathOperator{\Cov}{\mathbb{C}ov}
% \newcommand{\widesim}[2][1.5]{
%   \mathrel{\overset{#2}{\scalebox{#1}[1]{$\sim$}}}
% }
% \newcommand{\simiid}{\widesim{\phantom{\text{.}}\text{iid}\phantom{\text{.}}}}
% 
% \DeclareMathOperator{\Ber}{Ber}
% \DeclareMathOperator{\Bin}{Bin}
% \DeclareMathOperator{\Geom}{Geom}
% \DeclareMathOperator{\Poi}{Poi}
% \DeclareMathOperator{\Hyp}{Hyp}
% \DeclareMathOperator{\Unif}{\mathcal{U}}
% \DeclareMathOperator{\Exp}{Exp}
% \DeclareMathOperator{\Norm}{\mathcal{N}}
% \DeclareMathOperator{\Gam}{Gamma}



% SECTION: Theorem styles
\newtheoremstyle{break}
  {\topsep}{\topsep}%
  {\itshape}{}%
  {\bfseries}{}%
  {\newline}{}%

\newtheoremstyle{broke}
  {\topsep}{\topsep}%
  {}{}%
  {\bfseries}{}%
  {\newline}{}%

% \newtheoremstyle{definition}
%   {\topsep}{\topsep}%
%   {}{}%
%   {\bfseries}{}%
%   {\newline}{}%

% SECTION: Standard theorem blocks
\theoremstyle{broke}
% NOTE: Non-`\usepackage{mdframed}` definitions
% \newtheorem{theorem}{Theorem}[subsection]
% \newtheorem{proposition}[theorem]{Proposition}
% \newtheorem{lemma}[theorem]{Lemma}
% \newtheorem{corollary}[theorem]{Corollary}
% \newtheorem{claim}[theorem]{Claim}
% \theoremstyle{definition}
% \newtheorem{definition}[theorem]{Definition}

% NOTE: Definitions using `\usepackage{mdframed}`. 
\newmdtheoremenv[
  leftline = true,%
  rightline = false,%
  topline = false,%
  bottomline = false,%
  leftmargin = 0,%
  innertopmargin = 0,%
  innerbottommargin = 0,%
  innerrightmargin = 0,%
]{theorem}{Theorem}[subsection]
\newmdtheoremenv[
  leftline = true,%
  rightline = false,%
  topline = false,%
  bottomline = false,%
  leftmargin = 0,%
  innertopmargin = 0,%
  innerbottommargin = 0,%
  innerrightmargin = 0,%
]{proposition}[theorem]{Proposition}
\newmdtheoremenv[
  leftline = true,%
  rightline = false,%
  topline = false,%
  bottomline = false,%
  leftmargin = 0,%
  innertopmargin = 0,%
  innerbottommargin = 0,%
  innerrightmargin = 0,%
]{lemma}[theorem]{Lemma}
\newmdtheoremenv[
  leftline = true,%
  rightline = false,%
  topline = false,%
  bottomline = false,%
  leftmargin = 0,%
  innertopmargin = 0,%
  innerbottommargin = 0,%
  innerrightmargin = 0,%
]{corollary}[theorem]{Corollary}



% SECTION: QED symbol
\newcommand{\bsq}{\hfill$\blacksquare$}
% Current QED symbol is:
%   An amogus from \usepackage{tikz-among-us}
%   This may generate a warning each time `\qed` is called. Ignore it.
\renewcommand{\qedsymbol}{
  \begin{tikzpicture}
      \amongUsIII[xscale = -0.04, yscale = 0.04]{black}{white}
  \end{tikzpicture}
}

% SECTION: Goal management
\newcommand{\pf}{\pan\textbf{\underline{Proof.}}}
\newcommand{\sol}{\pan\textbf{Solution.}}
\newcommand{\pfthat}[1]{\pan\textbf{\underline{Proof that {#1}.}}}
\newcommand{\pfof}[1]{\pan\textbf{\underline{Proof of {#1}.}}}
\newcommand{\goal}[1]{
  \vspace{0.25\baselineskip}
  \pan\textbf{\underline{#1}}
  \vspace{0.25\baselineskip}
}
% \newcommand{\smallgoal}[1]{\par\textit{\underline{(#1):}}}
\newcommand{\case}[1]{
  \vspace{0.25\baselineskip}
  \pan\textit{\underline{(Case #1):}}
  \vspace{0.25\baselineskip}
}

\newmdenv[%
  leftline = true,%
  rightline = false,%
  topline = false,%
  bottomline = false,%
  leftmargin = 15,%
  innertopmargin = 0,%
  innerbottommargin = 0,%
  innerleftmargin = 15,%
  innerrightmargin = 0,%
]{leftish}
\newcommand{\withGoal}[2]{
  \vspace{0.25 \baselineskip}%
  \pan\textbf{\underline{#1}}%
  \vspace{-0.25 \baselineskip}%
  \begin{leftish}%
    #2%
  \end{leftish}%
}
\newcommand{\withCase}[2]{
  \vspace{0.25 \baselineskip}%
  \pan\textit{\underline{#1}}%
  \vspace{-0.25 \baselineskip}%
  \begin{leftish}%
    #2%
  \end{leftish}%
}



% SECTION: Minted setup
\setminted{linenos = true, breaklines = true, breakanywhere = true}



% SECTION: Bibliography style
% Labels \cite{Lee13} as [Lee13] in-text.
\DeclareFieldFormat{labelalpha}{\thefield{entrykey}}
\DeclareFieldFormat{extraalpha}{}



% SECTION: Obligatory
\allowdisplaybreaks
